\input{../../../templates/es/preambulo.tex}

% Los listados se toman de src/ con \lstinputlisting; la diapositiva es el archivo.
\lstset{
    basicstyle=\ttfamily\fontsize{8pt}{9.6pt}\selectfont,
    breaklines=true,
    breakatwhitespace=true,
    columns=fullflexible,
    keepspaces=true,
    xleftmargin=0.3em,
    framesep=2pt,
    aboveskip=0pt,
    belowskip=0pt
}

\newcommand{\marcoresultado}[3]{%
    \begin{frame}{#1}
        \begin{center}
            \includegraphics[height=0.62\textheight,width=\textwidth,keepaspectratio]{#2}
        \end{center}
        \vspace{0.25em}
        {\small #3\par}
    \end{frame}%
}

\newcommand{\marcoafirmacion}[2]{%
    \begin{frame}{#1}
        {\small #2\par}
    \end{frame}%
}

\date{}

\begin{document}

\tituloleccion{08}{Funciones de Caja Negra}

\section{Esta lección}

\begin{frame}{Esta lección}
La Lección 07 afinó perillas en un paisaje que todavía se veía. Esta función llega de otro lado: cinco argumentos, tres tipos, sin gráfica honesta.

\vspace{0.6em}
\begin{itemize}
    \item Una caja negra declara entradas legales y no dice nada de la salida.
    \item Genes reales, enteros y categóricos mezclados necesitan reglas de legalidad.
    \item Doce búsquedas que coinciden en \(x\) mientras la aptitud difiere por miles no son un máximo.
\end{itemize}
\end{frame}

% ================= caja_negra_01_sin_imagen =================
\begin{frame}{Un dominio mixto y un objetivo peligroso}
\[
\mathcal X=[a_l,a_u]\times[b_l,b_u]\times[x_l,x_u]
\times\mathcal N\times\mathcal F.
\]
Reparar proyecta propuestas a este producto cartesiano y cambia su distribución.
Aquí \(D=(n+1)^2(1+a+b)(120-x^2)r_\phi(x;b,n)+\tfrac12\) cambia de signo.
La bisección localiza \(D(x^*)=0\). Ese punto es un polo: el problema es no acotado o está mal
planteado en un dominio que contiene \(x^*\), no ``resuelto''.
\end{frame}

\section{La función que no puedes mirar}

\begin{frame}{La función que no puedes mirar}
Cada lección hasta ahora ha optimizado f(x) = sin(x) - 0.2*|x| o un pariente cercano:
un gen real, una curva, una imagen en la diapositiva y - crucialmente - un
óptimo por fuerza bruta en una cuadrícula para comprobar la ejecución. El Paso 5 de la Lección 02 lo dijo
con esas palabras y advirtió que la respuesta honesta a "¿tuvo éxito esta ejecución?"
se vuelve difícil una vez que x deja de ser un solo número.
\end{frame}

\marcoresultado{La función que no puedes mirar}{../figures/caja_negra_01_sin_imagen.png}{%
    Cinco cortes vivos coinciden en \texttt{x = -10.950}, pero sus máximos difieren por \texttt{9.23×}; la imagen no ha resuelto el valor}

% ================= caja_negra_02_la_cuadricula =================
\section{La cuadrícula, tasada y luego refinada}

\begin{frame}{La cuadrícula, tasada y luego refinada}
Cada lección anterior verificó sus ejecuciones contra un óptimo de fuerza bruta en una cuadrícula.
Esa comprobación es lo que hizo que "¿tuvo éxito esta ejecución?" fuera respondible. Este paso intenta
comprarlo aquí, y falla dos veces.
\end{frame}

\begin{frame}[fragile]{(1) evaluation\_cost()}
\lstinputlisting[basicstyle=\ttfamily\fontsize{7pt}{8.5pt}\selectfont,firstline=67,lastline=87]{../src/caja_negra_02_la_cuadricula.py}
\end{frame}

\begin{frame}[fragile]{(2) BOOK\_GRID}
\lstinputlisting[basicstyle=\ttfamily\fontsize{8pt}{9.6pt}\selectfont,firstline=89,lastline=101]{../src/caja_negra_02_la_cuadricula.py}
\end{frame}

\begin{frame}[fragile]{(3) grid\_maximum()}
\lstinputlisting[basicstyle=\ttfamily\fontsize{7pt}{8.5pt}\selectfont,firstline=120,lastline=143]{../src/caja_negra_02_la_cuadricula.py}
\end{frame}

\begin{frame}[fragile]{(4) REFINEMENTS}
\lstinputlisting[basicstyle=\ttfamily\fontsize{8pt}{9.6pt}\selectfont,firstline=145,lastline=154]{../src/caja_negra_02_la_cuadricula.py}
\end{frame}

\marcoresultado{La cuadrícula, tasada y luego refinada}{../figures/caja_negra_02_la_cuadricula.png}{%
    La cuadrícula del libro cuesta 105,525,000 llamadas; el refinamiento hace que el máximo reportado crezca 61,099 veces en lugar de converger}

% ================= caja_negra_03_el_cromosoma =================
\section{Un cromosoma, tres tipos de gen}

\begin{frame}{Un cromosoma, tres tipos de gen}
Hasta ahora un cromosoma era un número real, o un vector de números reales, y
"legal" significaba "dentro de [-10, 10]". Aquí un individuo porta tres tipos
diferentes de gen a la vez:
\end{frame}

\begin{frame}[fragile]{(1) clamp()}
\lstinputlisting[basicstyle=\ttfamily\fontsize{8pt}{9.6pt}\selectfont,firstline=80,lastline=89]{../src/caja_negra_03_el_cromosoma.py}
\end{frame}

\begin{frame}[fragile]{(2) closest()}
\lstinputlisting[basicstyle=\ttfamily\fontsize{8pt}{9.6pt}\selectfont,firstline=91,lastline=102]{../src/caja_negra_03_el_cromosoma.py}
\end{frame}

\begin{frame}[fragile]{(3) Individual}
\lstinputlisting[basicstyle=\ttfamily\fontsize{6pt}{7.2pt}\selectfont,firstline=104,lastline=134]{../src/caja_negra_03_el_cromosoma.py}
\end{frame}

\begin{frame}[fragile]{(4) create\_random()}
\lstinputlisting[basicstyle=\ttfamily\fontsize{8pt}{9.6pt}\selectfont,firstline=136,lastline=153]{../src/caja_negra_03_el_cromosoma.py}
\end{frame}

\begin{frame}[fragile]{(5) repair\_report()}
\lstinputlisting[basicstyle=\ttfamily\fontsize{8pt}{9.6pt}\selectfont,firstline=155,lastline=171]{../src/caja_negra_03_el_cromosoma.py}
\end{frame}

\marcoresultado{Un cromosoma, tres tipos de gen}{../figures/caja_negra_03_el_cromosoma.png}{%
    El constructor repara los cinco ejemplos ilegales; la población inicial muestrea solo el 10\% del rango declarado de \texttt{x}}

% ================= caja_negra_04_los_operadores =================
\section{Los operadores, y lo que cuesta la reparación silenciosa}

\begin{frame}{Los operadores, y lo que cuesta la reparación silenciosa}
El paso 3 puso toda la legalidad en el constructor, por lo que los operadores pueden escribirse
como si cada gen fuera un número real. Este paso los escribe - el diseño del libro,
y uno deliberado: cada gen tiene su propia moneda y su propio tamaño de paso,
porque los genes son diferentes tipos de objeto. Mezcla alfa 0.3 para
a, 0.5 para b, 1.0 para x; sigma gaussiano 0.2 para a y b, 1.0 para x y n;
el gen etiqueta muta volviéndose a elegir, la única mutación que admite un conjunto.
\end{frame}

\begin{frame}[fragile]{(1) cruza()}
\lstinputlisting[basicstyle=\ttfamily\fontsize{7pt}{8.5pt}\selectfont,firstline=126,lastline=151]{../src/caja_negra_04_los_operadores.py}
\end{frame}

\begin{frame}[fragile]{(2) mutar()}
\lstinputlisting[basicstyle=\ttfamily\fontsize{7pt}{8.5pt}\selectfont,firstline=153,lastline=174]{../src/caja_negra_04_los_operadores.py}
\end{frame}

\begin{frame}[fragile]{(3) el censo de rep.}
\lstinputlisting[basicstyle=\ttfamily\fontsize{6pt}{7.2pt}\selectfont,firstline=192,lastline=223]{../src/caja_negra_04_los_operadores.py}
\end{frame}

\begin{frame}[fragile]{(4) la demo de deriva}
\lstinputlisting[basicstyle=\ttfamily\fontsize{7pt}{8.5pt}\selectfont,firstline=225,lastline=248]{../src/caja_negra_04_los_operadores.py}
\end{frame}

\marcoresultado{Los operadores, y lo que cuesta la reparación silenciosa}{../figures/caja_negra_04_los_operadores.png}{%
    41.5\% de las mutaciones enteras forzadas vuelven a su estado original sin cambios; mallas gruesas pueden silenciar la mutación por completo}

% ================= caja_negra_05_la_busqueda =================
\section{La búsqueda, y una respuesta que no se queda quieta}



\begin{frame}[fragile]{(1) select\_rank\_with\_elite()}
\lstinputlisting[basicstyle=\ttfamily\fontsize{7pt}{8.5pt}\selectfont,firstline=166,lastline=190]{../src/caja_negra_05_la_busqueda.py}
\end{frame}

\begin{frame}[fragile]{(2) el ciclo generacional}
\lstinputlisting[basicstyle=\ttfamily\fontsize{7pt}{8.5pt}\selectfont,firstline=193,lastline=218]{../src/caja_negra_05_la_busqueda.py}
\end{frame}

\begin{frame}[fragile]{(3) los diagnósticos}
\lstinputlisting[basicstyle=\ttfamily\fontsize{7pt}{8.5pt}\selectfont,firstline=228,lastline=253]{../src/caja_negra_05_la_busqueda.py}
\end{frame}

\marcoresultado{La búsqueda, y una respuesta que no se queda quieta}{../figures/caja_negra_05_la_busqueda.png}{%
    La mejor aptitud alcanza \texttt{2.278e-09}, pero la curva sigue subiendo y el 97.5\% de la población final ha salido del intervalo inicial de \texttt{x}}

% ================= caja_negra_06_el_veredicto =================
\section{El veredicto, doce veces más}

\begin{frame}{El veredicto, doce veces más}
La ejecución del paso 5 pareció un éxito, y una ejecución es una anécdota. Este paso
repite la búsqueda completa doce veces, desde doce semillas, y pone a los
doce campeones lado a lado. Nada sobre el algoritmo cambia: mismo
cromosoma, mismos operadores, mismas perillas, mismas 100 generaciones. Solo la semilla
se mueve.
\end{frame}

\begin{frame}[fragile]{(1) run\_search()}
\lstinputlisting[basicstyle=\ttfamily\fontsize{7pt}{8.5pt}\selectfont,firstline=180,lastline=203]{../src/caja_negra_06_el_veredicto.py}
\end{frame}

\begin{frame}[fragile]{(2) el Monte Carlo}
\lstinputlisting[basicstyle=\ttfamily\fontsize{8pt}{9.6pt}\selectfont,firstline=208,lastline=225]{../src/caja_negra_06_el_veredicto.py}
\end{frame}

\begin{frame}[fragile]{(3) las dos lecturas}
\lstinputlisting[basicstyle=\ttfamily\fontsize{8pt}{9.6pt}\selectfont,firstline=227,lastline=243]{../src/caja_negra_06_el_veredicto.py}
\end{frame}

\marcoresultado{El veredicto, doce veces más}{../figures/caja_negra_06_el_veredicto.png}{%
    Los campeones ocupan un intervalo estrecho de \texttt{x} mientras que su aptitud difiere en 21,504×, lo cual es inconsistente con la convergencia a un máximo finito}

% ================= caja_negra_07_abriendo_la_caja =================
\section{Abriendo la caja}

\begin{frame}{Abriendo la caja}
Seis pasos han tratado complicated\_one() como una caja negra, la manera en que el libro
la presenta. Las búsquedas están hechas; ahora leemos la única línea que estaba haciendo
todo el daño. Cada término de la suma es un numerador sobre un denominador, y
el denominador es
\end{frame}

\begin{frame}[fragile]{(1) denominator()}
\lstinputlisting[basicstyle=\ttfamily\fontsize{8pt}{9.6pt}\selectfont,firstline=75,lastline=87]{../src/caja_negra_07_abriendo_la_caja.py}
\end{frame}

\begin{frame}[fragile]{(2) el polo}
\lstinputlisting[basicstyle=\ttfamily\fontsize{7pt}{8.5pt}\selectfont,firstline=98,lastline=119]{../src/caja_negra_07_abriendo_la_caja.py}
\end{frame}

\begin{frame}[fragile]{(3) la escalera}
\lstinputlisting[basicstyle=\ttfamily\fontsize{7pt}{8.5pt}\selectfont,firstline=121,lastline=141]{../src/caja_negra_07_abriendo_la_caja.py}
\end{frame}

\marcoresultado{Abriendo la caja}{../figures/caja_negra_07_abriendo_la_caja.png}{%
    El óptimo aparente es un polo cerca de \texttt{x = -10.9535836}; la aptitud reportada mide la distancia a la singularidad}

\section{Siguiente}

\begin{frame}{Siguiente}
El óptimo aparente es un polo cerca de \(x = -10.9535836\). La aptitud reportada era distancia a una singularidad.

\vspace{0.8em}
\textbf{Lección 09 --- Optimización combinatoria binaria}

La imagen se fue para siempre. Los bits codifican selección, asignación, colocación. La factibilidad tiene que seguir visible: penalización, reparación, o una comprobación exacta.
\end{frame}
\end{document}
